% !TEX TS-program = XeLaTeX
% use the following command:
% all document files must be coded in UTF-8
\documentclass[spanish]{textolivre}
% build HTML with: make4ht -e build.lua -c textolivre.cfg -x -u article "fn-in,svg,pic-align"

\journalname{Texto Livre}
\thevolume{19}
%\thenumber{1} % old template
\theyear{2026}
\receiveddate{\DTMdisplaydate{2025}{12}{17}{-1}} % YYYY MM DD
\accepteddate{\DTMdisplaydate{2026}{2}{28}{-1}}
\publisheddate{\DTMdisplaydate{2026}{9}{1}{-1}}
\corrauthor{Gloria Mónica Martínez Aguilar}
\articledoi{10.1590/1983-3652.2026.63539}
%\articleid{NNNN} % if the article ID is not the last 5 numbers of its DOI, provide it using \articleid{} commmand 
% list of available sesscions in the journal: articles, dossier, reports, essays, reviews, interviews, editorial
\articlesessionname{articles}
\runningauthor{Martínez Aguilar, Segura Félix y Guerrero López} 
%\editorname{Leonardo Araújo} % old template
\sectioneditorname{Hugo Heredia Ponce~\orcid{0000-0003-3657-1369}}
\layouteditorname{Rayany Chaves~\orcid{0009-0007-0028-484X}}

\title{Percepción estudiantil de la inteligencia artificial generativa en procesos de estudio: un análisis del uso de ChatGPT en educación superior}
\othertitle{Percepção estudantil da inteligência artificial generativa nos processos de estudo: uma análise do uso do ChatGPT no ensino superior}
\othertitle{Student Perception of Generative Artificial Intelligence in Study Processes: An Analysis of ChatGPT Use in Higher Education}
% if there is a third language title, add here:
%\othertitle{Artikelvorlage zur Einreichung beim Texto Livre Journal}

\author[1]{Gloria Mónica Martínez Aguilar~\orcid{0000-0003-3834-4880}\thanks{Email: \href{mailto:ketherelohim@gmail.com}{ketherelohim@gmail.com}}}
\author[1]{Kristian Segura Félix~\orcid{0000-0002-0820-5620}\thanks{Email: \href{mailto:ksegura984@gmail.com}{ksegura984@gmail.com}}}
\author[2]{Geshel David Guerrero López~\orcid{0000-0003-0341-501X}\thanks{Email: \href{mailto:gesheldavid.gl@nlaredo.tecnm.mx}{gesheldavid.gl@nlaredo.tecnm.mx}}}
\affil[1]{Universidad Tecnológica de Torreón, Torreón, Coahuila, México.}
\affil[2]{Instituto Tecnológico de Nuevo Laredo, Departamento de Eléctrica Electrónica, Nuevo Laredo, Tamaulipas, México.}

\addbibresource{article.bib}
% use biber instead of bibtex
% $ biber article

% used to create dummy text for the template file
\definecolor{dark-gray}{gray}{0.35} % color used to display dummy texts
\usepackage{lipsum}
\SetLipsumParListSurrounders{\colorlet{oldcolor}{.}\color{dark-gray}}{\color{oldcolor}}

% used here only to provide the XeLaTeX and BibTeX logos
\usepackage{hologo}

% if you use multirows in a table, include the multirow package
\usepackage{multirow}

% provides sidewaysfigure environment
\usepackage{rotating}

% CUSTOM EPIGRAPH - BEGIN 
%%% https://tex.stackexchange.com/questions/193178/specific-epigraph-style
\usepackage{epigraph}
\renewcommand\textflush{flushright}
\makeatletter
\newlength\epitextskip
\pretocmd{\@epitext}{\em}{}{}
\apptocmd{\@epitext}{\em}{}{}
\patchcmd{\epigraph}{\@epitext{#1}\\}{\@epitext{#1}\\[\epitextskip]}{}{}
\makeatother
\setlength\epigraphrule{0pt}
\setlength\epitextskip{0.5ex}
\setlength\epigraphwidth{.7\textwidth}
% CUSTOM EPIGRAPH - END

% to use IPA symbols in unicode add
%\usepackage{fontspec}
%\newfontfamily\ipafont{CMU Serif}
%\newcommand{\ipa}[1]{{\ipafont #1}}
% and in the text you may use the \ipa{...} command passing the symbols in unicode

% LANGUAGE - BEGIN
% ARABIC
% for languages that use special fonts, you must provide the typeface that will be used
% \setotherlanguage{arabic}
% \newfontfamily\arabicfont[Script=Arabic]{Amiri}
% \newfontfamily\arabicfontsf[Script=Arabic]{Amiri}
% \newfontfamily\arabicfonttt[Script=Arabic]{Amiri}
%
% in the article, to add arabic text use: \textlang{arabic}{ ... }
%
% RUSSIAN
% for russian text we also need to define fonts with support for Cyrillic script
% \usepackage{fontspec}
% \setotherlanguage{russian}
% \newfontfamily\cyrillicfont{Times New Roman}
% \newfontfamily\cyrillicfontsf{Times New Roman}[Script=Cyrillic]
% \newfontfamily\cyrillicfonttt{Times New Roman}[Script=Cyrillic]
%
% in the text use \begin{russian} ... \end{russian}
% LANGUAGE - END

% EMOJIS - BEGIN
% to use emoticons in your manuscript
% https://stackoverflow.com/questions/190145/how-to-insert-emoticons-in-latex/57076064
% using font Symbola, which has full support
% the font may be downloaded at:
% https://dn-works.com/ufas/
% add to preamble:
% \newfontfamily\Symbola{Symbola}
% in the text use:
% {\Symbola }
% EMOJIS - END

% LABEL REFERENCE TO DESCRIPTIVE LIST - BEGIN
% reference itens in a descriptive list using their labels instead of numbers
% insert the code below in the preambule:
%\makeatletter
%\let\orgdescriptionlabel\descriptionlabel
%\renewcommand*{\descriptionlabel}[1]{%
%  \let\orglabel\label
%  \let\label\@gobble
%  \phantomsection
%  \edef\@currentlabel{#1\unskip}%
%  \let\label\orglabel
%  \orgdescriptionlabel{#1}%
%}
%\makeatother
%
% in your document, use as illustraded here:
%\begin{description}
%  \item[first\label{itm1}] this is only an example;
%  % ...  add more items
%\end{description}
% LABEL REFERENCE TO DESCRIPTIVE LIST - END


% add line numbers for submission
%\usepackage{lineno}
%\linenumbers

\begin{document}
\maketitle

\begin{polyabstract}
\begin{abstract}
El uso de herramientas basadas en inteligencia artificial generativa, como ChatGPT, se ha expandido rápidamente en la educación superior, generando oportunidades de apoyo académico, pero también inquietudes relacionadas con la autonomía, la ética y la integridad estudiantil. Este estudio analizó las percepciones de 382 estudiantes universitarios de licenciatura, ingeniería y posgrado respecto al uso académico de ChatGPT. Se aplicó un cuestionario Likert de 12 ítems organizado en siete dimensiones: utilidad percibida, autonomía, riesgos, competencias cognitivas, confianza, políticas institucionales y criterio ético. Los resultados descriptivos mostraron una valoración moderadamente favorable hacia el uso académico de la herramienta. El análisis inferencial indicó ausencia de diferencias significativas por género, mientras que el nivel de preparación académica presentó diferencias únicamente en la dimensión de autonomía, donde los estudiantes intermedios y avanzados reportaron mayor confianza para emplear ChatGPT en actividades académicas. El análisis por edad tampoco evidenció diferencias significativas entre grupos. Estos hallazgos sugieren que la integración de IA generativa es vista de manera relativamente homogénea entre perfiles estudiantiles diversos, aunque persisten necesidades de orientación institucional y formación ética para su uso responsable en contextos educativos.

\keywords{Inteligencia artificial \sep Educación superior \sep Percepciones estudiantiles \sep ChatGPT}
\end{abstract}

\begin{portuguese}
\begin{abstract}
O uso de ferramentas baseadas em inteligência artificial generativa, como o ChatGPT, expandiu-se rapidamente no ensino superior, gerando oportunidades de apoio acadêmico, mas também levantando preocupações relacionadas à autonomia, ética e integridade estudantil. Este estudo analisou as percepções de 382 estudantes de graduação, Engenharia e pós-graduação sobre o uso acadêmico do ChatGPT. Foi aplicado um questionário com escala Likert de 12 itens, organizado em sete dimensões: utilidade percebida, autonomia, riscos, habilidades cognitivas, confiança, políticas institucionais e considerações éticas. Os resultados descritivos mostraram uma avaliação moderadamente favorável do uso acadêmico da ferramenta. A análise inferencial não indicou diferenças significativas por gênero, enquanto o nível de preparação acadêmica apresentou diferenças apenas na dimensão autonomia, em que estudantes de nível intermediário e avançado relataram maior confiança no uso do ChatGPT para atividades acadêmicas. A análise por idade também não revelou diferenças significativas entre os grupos. Esses achados sugerem que a integração da IA generativa é vista de forma relativamente homogênea entre os diversos perfis de estudantes, embora haja uma necessidade persistente de orientação institucional e treinamento ético para seu uso responsável em contextos educacionais.

\keywords{Inteligência artificial \sep Ensino superior \sep Percepções dos alunos \sep ChatGPT}
\end{abstract}
\end{portuguese}

\begin{english}
\begin{abstract}
The use of generative artificial intelligence tools, such as ChatGPT, has rapidly expanded across higher education, offering new forms of academic support while raising concerns related to student autonomy, ethics, and academic integrity. This study examined the perceptions of 382 undergraduate, engineering, and graduate students regarding the academic use of ChatGPT. A 12-item Likert questionnaire was administered, structured into seven dimensions: perceived usefulness, autonomy, perceived risks, cognitive competencies, trust, institutional policies, and ethical criteria. Descriptive results indicated a moderately favorable evaluation of the tool. Inferential analyses showed no significant differences by gender, whereas academic preparation level revealed differences only in the autonomy dimension, with intermediate and advanced students reporting greater confidence in using ChatGPT for academic tasks. Age-based comparisons also showed no significant differences between groups. Overall, findings suggest that generative AI integration is perceived relatively consistently across diverse student profiles, although there remains a need for institutional guidance and ethical training to promote responsible use in educational contexts.

\keywords{Artificial intelligence \sep Higher education \sep Student perceptions \sep ChatGPT}
\end{abstract}
\end{english}
% if there is another abstract, insert it here using the same scheme
\end{polyabstract}

\section{ Introducción}\label{sec-intro}

%Rever o título da seção, porque o artigo é em espanhol

El avance reciente de la inteligencia artificial generativa (IAG) ha transformado aceleradamente los procesos educativos en la educación superior. Herramientas como ChatGPT se han incorporado a actividades de estudio, resolución de problemas, elaboración de textos, programación y apoyo a la investigación, modificando tanto las prácticas de aprendizaje como las dinámicas institucionales. Diversos estudios señalan que estas tecnologías pueden ampliar las oportunidades para el aprendizaje autónomo, ofrecer retroalimentación instantánea a gran escala, especialmente por su capacidad para procesar lenguaje natural, generar explicaciones personalizadas, facilitar la comprensión de contenidos y asistir en la resolución de tareas complejas \cite{BaidooAnuOwusuAnsah2023,CromptonBurke2023,Dempere2023,MaiVanDaVanHanh2024,ZawackiRichterMarinBondGouverneur2019}.

No obstante, su expansión ha generado preocupaciones en torno a la integridad académica, la originalidad de los trabajos y la precisión de los contenidos generados \cite{Cotton2024,Dwivedi2023,Luo2024}.Ante este escenario, diversas instituciones de educación superior alrededor del mundo han comenzado a revisar sus marcos normativos, prácticas de evaluación y lineamientos de uso, lo que subraya la necesidad de comprender cómo el estudiantado interpreta e incorpora estas herramientas en sus procesos formativos \cite{MichelVillarreal2023,Zhang2025}.

En este contexto, la percepción estudiantil constituye un indicador clave para analizar la integración efectiva de la IAG. Investigaciones recientes reportan variaciones asociadas a disciplina, nivel académico, género y experiencia previa con inteligencia artificial (IA) \cite{BlahopoulouOrtizBonnin2025,KlimovaDeCampos2024,StohrOuMalmstrom2024}.Esto sugiere que la adopción de ChatGPT no responde a un patrón homogéneo, sino a condiciones formativas y contextuales diferenciadas \cite{Ravselj2025}.

A pesar de estos avances, la evidencia disponible se concentra principalmente en Europa, Norteamérica y Asia, lo que abre una brecha de conocimiento sobre las percepciones de estudiantes latinoamericanos, especialmente en instituciones tecnológicas y programas de Ciencia, Tecnología, Ingeniería y Matemáticas (STEM, por sus siglas en inglés) y dificulta el diseño de lineamientos contextualizados y estrategias docentes sustentadas en las experiencias reales del estudiantado.

En respuesta a esta necesidad, el presente estudio analiza las percepciones de estudiantes universitarios sobre el uso académico de ChatGPT, explorando siete dimensiones conceptuales: utilidad, autonomía, riesgos, competencias cognitivas, confianza, políticas institucionales y criterio ético. Asimismo, compara dichas percepciones según género, nivel de preparación académica y categoría de edad, con el fin de identificar patrones diferenciales y aportar evidencia para la construcción de lineamientos institucionales orientados a un uso responsable y pedagógicamente significativo de la IAG.

\section{Marco Teórico}

\subsection{Inteligencia artificial generativa en educación superior}
Herramientas como ChatGPT permiten automatizar procesos de revisión, análisis textual y síntesis de información, lo que ha favorecido su adopción en prácticas académicas tanto del estudiantado como del profesorado. Sin embargo, esto plantea discusiones relacionadas con su impacto en las competencias cognitivas, la autonomía del aprendizaje y la dependencia de sistemas automatizados \cite{Kasneci2023,RudolphTanTan2023}. Las discusiones actuales giran en torno a cómo equilibrar el uso pedagógico de estas herramientas con el desarrollo de pensamiento crítico, creatividad y habilidades de resolución de problemas.

\subsection{Beneficios educativos y oportunidades de aprendizaje}

Pese a las preocupaciones por el uso de herramientas como ChatGPT, investigaciones destacan su potencial para fortalecer la autonomía del estudiantado al permitir consultas ilimitadas y respuestas adaptadas al nivel del usuario, lo que puede favorecer estrategias metacognitivas \cite{Dempere2023}.

De manera complementaria, otras investigaciones señalan ventajas relacionadas con la reducción de barreras idiomáticas, el apoyo en traducción automatizada y el desarrollo de habilidades de escritura, contribuyendo así a entornos de aprendizaje más inclusivos \cite{WangXuLiu2024}. En este sentido, el valor pedagógico de ChatGPT radica en su integración como herramienta guiada dentro de actividades formativas con propósitos definidos.

\subsection{Riesgos, desafíos y preocupaciones éticas}

No obstante, este potencial educativo debe examinarse de manera crítica, desde esta perspectiva la literatura también advierte riesgos relevantes, como la posible sustitución del esfuerzo cognitivo, la dificultad para distinguir contenido generado por IA y la proliferación de prácticas deshonestas en entornos de evaluación \cite{Cotton2024,Lee2024}. Asimismo, se reportan problemas relacionados con la precisión variable de las respuestas, la presencia de sesgos en los modelos de lenguaje y la reproducción automática de información incorrecta o no verificada \cite{Dwivedi2023,Kasneci2023}.

En el plano institucional, \textcite{Luo2024} argumenta que la noción tradicional de originalidad debe ser revisada en la era de la IAG, mientras que \textcite{BlahopoulouOrtizBonnin2025} muestran que estudiantes usuarios y no usuarios difieren en la forma de percibir estos riesgos. En conjunto, estos hallazgos muestran que la integración educativa de ChatGPT no solo plantea retos técnicos, sino también pedagógicos, normativos y éticos.

\subsection{Percepciones estudiantiles sobre ChatGPT}

La evidencia empírica indica que las percepciones estudiantiles son multidimensionales y están influidas por factores como la disciplina, el nivel académico, el género y la familiaridad previa con IA \cite{KlimovaDeCampos2024,StohrOuMalmstrom2024}. Estudios como el de \textcite{Ravselj2025} muestran un patrón consistente: el estudiantado reconoce el valor pedagógico de ChatGPT, pero al mismo tiempo expresa preocupaciones relacionadas con la ética, la privacidad, la dependencia y la fiabilidad de las respuestas.

Estas diferencias entre grupos refuerzan la necesidad de examinar el fenómeno en contextos educativos específicos. Desde esta perspectiva, el presente estudio aporta evidencia sobre un entorno poco explorado en la literatura, al analizar cuantitativamente las percepciones del estudiantado universitario latinoamericano y las diferencias observadas según variables académicas y sociodemográficas.


\section{Metodología}

\subsection{Enfoque y diseño del estudio}
Este estudio adoptó un diseño cuantitativo, no experimental, transversal y de alcance descriptivo–comparativo, cuyo propósito fue analizar las percepciones estudiantiles sobre el uso académico de ChatGPT en educación superior. El estudio se centró en identificar niveles de aceptación, riesgos percibidos, autonomía en el uso de IA y consideraciones éticas, así como en examinar diferencias según género, nivel de formación académica y categoría de edad.


\subsection{Participantes}

La población de referencia estuvo conformada por aproximadamente 4000 estudiantes de la Universidad Tecnológica de Torreón. La muestra estuvo integrada por 382 estudiantes universitarios inscritos en programas de licenciatura, ingeniería y posgrado, seleccionados mediante un muestreo no probabilístico por conveniencia. Esta composición permitió contar con una muestra heterogénea de distintos campos de formación. En el nivel de posgrado participaron estudiantes de la Maestría en Ingeniería en Manufactura Inteligente (MIMI) y de la Maestría en Economía Circular (MEC). En el nivel de ingeniería estuvieron presentes las carreras de Ingeniería en Mantenimiento Industrial, Ingeniería Mecánica, Ingeniería Mecatrónica, Ingeniería Industrial e Ingeniería en Tecnologías de la Información e Innovación Digital. Asimismo, se contó con la participación de estudiantes de la Licenciatura en Administración, lo que permitió integrar perspectivas tanto técnicas como administrativas y enriquecer el análisis.

La recolección de la información se realizó en el periodo del 19 al 30 de septiembre de 2025. La participación fue voluntaria y anónima. Se invitó a participar a 382 alumnos y todos respondieron el cuestionario, por lo que se trabajó con 382 casos válidos. No se registraron datos faltantes en la base analizada. Se recabaron variables sociodemográficas como género, edad y cuatrimestre. En la muestra participaron 126 mujeres y 256 hombres. El cuatrimestre cursado se transformó en tres niveles de preparación académica: Básico (1–6), con 245 alumnos; Intermedio (7–10), con 122 alumnos; y Avanzado (11–16), con 15 alumnos. De manera complementaria, la edad se recategorizó en tres grupos: Jóvenes (17–21 años), con 242 alumnos; Adultos emergentes (22–29 años), con 99 alumnos; y Adultos profesionales (30–50 años), con 41 alumnos. Estas transformaciones permitieron realizar análisis comparativos entre diferentes etapas formativas, rangos de desarrollo y género.

\subsection{Instrumento}

Se aplicó un cuestionario elaborado en Google Forms, conformado por 12 ítems tipo Likert de 5 puntos (1 = “Totalmente en desacuerdo”, 5 = “Totalmente de acuerdo”), diseñado para evaluar percepciones estudiantiles sobre el uso académico de ChatGPT. Los ítems se organizaron en siete dimensiones teóricas, construidas a partir de literatura reciente sobre IA aplicada a contextos educativos:

\begin{itemize}
    \item Utilidad percibida (UTI)
    \item Autonomía en el aprendizaje (AUT)
    \item Riesgos y ética académica (RIE)
    \item Competencias cognitivas (COM)
    \item Confianza en la precisión de las respuestas (CON)
    \item Políticas institucionales (POL)
    \item Criterio ético personal (CRI)
\end{itemize}
    

Para el análisis estadístico por dimensiones, se construyeron puntajes compuestos a partir del promedio de los reactivos correspondientes cuando la dimensión incluyó más de un ítem. De esta manera, la dimensión de utilidad percibida se integró con UTI1, UTI2 y UTI3; la dimensión de riesgos y ética académica con RIE1, RIE2 y RIE3; y la dimensión de políticas institucionales con POL1 y POL2. En los casos de autonomía en el aprendizaje, competencias cognitivas, confianza en la precisión de las respuestas y criterio ético personal, la medición se realizó a partir de un reactivo único para cada dimensión.

El cuestionario se administró de manera presencial durante sesiones de clase y mediante código QR. Se informó a las y los estudiantes sobre los objetivos de la investigación y se garantizó la confidencialidad de la información. No se recabaron datos sensibles, y el proceso se alineó con los lineamientos éticos universitarios.

\subsection{Análisis de datos}

Los análisis se realizaron en el software IBM SPSS Statistics 21. Las técnicas empleadas incluyeron:
\begin{itemize}
    \item Estadísticos descriptivos (media y desviación estándar) para los 12 ítems del cuestionario con el propósito de identificar tendencias generales de respuesta, nivel promedio de acuerdo y variabilidad de las percepciones estudiantiles. La media permitió establecer la dirección general de la valoración de cada reactivo, mientras que la desviación estándar permitió examinar el grado de dispersión de las respuestas en la muestra.
    \item Fiabilidad interna mediante alfa de Cronbach. Además del valor global del alfa, se revisaron las correlaciones ítem–total corregidas y el valor de alfa si se elimina el elemento, con el fin de estimar la contribución relativa de cada reactivo al comportamiento general del instrumento. Este análisis permitió valorar la coherencia interna de la escala en su conjunto y contar con un criterio adicional para sustentar la interpretación posterior de los resultados.
    \item Correlaciones de Pearson para identificar relaciones entre dimensiones del instrumento, con el objetivo de identificar la dirección e intensidad de las asociaciones lineales entre los distintos factores del instrumento. Para su interpretación se consideró tanto la magnitud del coeficiente como su significancia estadística.
    \item ANOVA de un factor para comparar las medias de las dimensiones según género, nivel de preparación académica y categoría de edad. 
    \item Pruebas post hoc Tukey. En los contrastes con más de dos grupos, cuando se identificaron diferencias estadísticamente significativas, se recurrió a pruebas post hoc de Tukey para precisar entre qué grupos se localizaban dichas diferencias. Antes de interpretar los ANOVA se revisó el supuesto de homogeneidad de varianzas mediante la prueba de Levene. En todos los análisis se trabajó con un nivel de significancia de p < .05 y con los 382 cuestionarios completos, sin exclusión de casos por datos faltantes.
\end{itemize}

    

El análisis se orientó a identificar patrones de percepción, diferencias significativas entre grupos y posibles implicaciones para promover un uso responsable de herramientas de IA en educación superior.

\section{Resultados}
\subsection{Estadísticos descriptivos de los ítems}

Los estadísticos descriptivos de los ítems, presentados en la Tabla \ref{tab1}, permiten identificar patrones generales en la percepción estudiantil sobre el uso académico de ChatGPT. Los valores medios se situaron entre 2.71 y 3.78, reflejando una actitud moderadamente favorable hacia la herramienta. Las desviaciones estándar, que oscilaron entre 1.03 y 1.39, evidencian una variabilidad amplia en las respuestas, lo que indica que, aunque existe una valoración general positiva, las percepciones no fueron homogéneas en todos los participantes.


\begin{table}[htbp]
\centering
\begin{threeparttable}
\caption{Estadísticos descriptivos de los ítems del cuestionario.}
\label{tab1}
\begin{tabular}{p{1cm} p{0.7cm} p{0.7cm} p{0.7cm} p{0.7cm} p{0.7cm} p{0.7cm} p{0.7cm} p{0.7cm} p{0.7cm} p{0.7cm} p{0.7cm} p{0.7cm}}
\toprule
Item & UTI1 & UTI2 & UTI3 & AUT1 & RIE1 & RIE2 & COM1 & CON1 & POL1 & POL2 & CRI1 & RIE3 \\ \midrule
Media & 3.78 & 3.54 & 3.41 & 3.21 & 3.75 & 3.20 & 3.68 & 3.01 & 3.16 & 2.71 & 3.41 & 2.92 \\
Desv. típ. & 1.166 & 1.207 & 1.278 & 1.290 & 1.356 & 1.294 & 1.212 & 1.031 & 1.276 & 1.397 & 1.173 & 1.279 \\ \bottomrule
\end{tabular}
\source{Estadísticos calculados por los autores en SPSS con datos del estudio. N=382.}
\end{threeparttable}
\end{table}

Los puntajes más altos se observaron en UTI1 (M = 3.78) y RIE1 (M = 3.75), lo que indica que los estudiantes reconocen la utilidad de ChatGPT, aunque también perciben riesgos asociados a su uso intensivo. Por otro lado, los valores más bajos correspondieron a POL2 (M = 2.71) y RIE3 (M = 2.92), lo que sugiere una percepción limitada sobre la existencia de lineamientos institucionales y una preocupación menos marcada en uno de los reactivos vinculados con riesgos específicos. En términos generales, los resultados descriptivos muestran que la aceptación de la herramienta convive con reservas relacionadas con su regulación y con algunos aspectos de su uso académico.


\subsection{Consistencia interna del instrumento}

La escala mostró una consistencia interna adecuada ($\alpha = .793$; $\alpha$ estandarizado = .801), calculada sobre los 12 reactivos del cuestionario. Estos valores indican una fiabilidad aceptable de la escala total para fines descriptivos y comparativos. Considerando que el instrumento integra dimensiones conceptualmente distintas, el coeficiente obtenido puede interpretarse como evidencia de una coherencia interna suficiente del conjunto de reactivos, sin implicar necesariamente un comportamiento unidimensional estricto de toda la escala.

Las correlaciones ítem–total fueron mayores en los reactivos asociados a utilidad percibida, autonomía, pensamiento crítico y comunicación (r = .49 – .69), indicando una contribución sólida a la medición global. En contraste, los ítems relacionados con riesgos percibidos y lineamientos institucionales mostraron correlaciones más bajas (r = .06 – .33), lo cual es congruente con el carácter multidimensional del instrumento y sustenta la interpretación diferenciada por dimensiones en los análisis posteriores. En este sentido, los resultados de fiabilidad permiten sostener el uso analítico del cuestionario, aunque también muestran que algunos reactivos se integran con menor fuerza a la escala total, especialmente aquellos vinculados con riesgos y regulación institucional.

\subsection{Correlaciones entre las dimensiones}

La Tabla \ref{tab2} presenta las correlaciones entre los factores del instrumento. Se identificó un conjunto de asociaciones fuertes entre utilidad percibida, autonomía, competencias cognitivas y criterio ético (r = .50 – .78, p < .01), lo que indica que las percepciones favorables hacia ChatGPT tienden a agruparse: estudiantes que encuentran útil la herramienta también reportan mayor autonomía y una valoración positiva de sus aportes en procesos académicos.


\begin{table}[htbp]
\centering
\begin{threeparttable}
\caption{Matriz de correlaciones de Pearson entre las dimensiones evaluadas.}
\label{tab2}
\begin{tabular}{lcccccc}
\toprule
Factor & UTI & RIE & POL & AUT & COM & CON \\ 
\midrule
UTI & — & & & & & \\
RIE & .045 & — & & & & \\
POL & .298 & .229 & — & & & \\
AUT & .778 & .024 & .258 & — & & \\
COM & .589 & .115* & .272 & .506** & — & \\
CON & .546 & -.011 & .294 & .487 & .363 & — \\
CRI & .556 & .078 & .304 & .477 & .486 & .378** \\
\bottomrule
\end{tabular}
\source{Datos del estudio.}
\notes{*p < .05; **p < .01.}
\end{threeparttable}
\end{table}

La dimensión de riesgos percibidos mostró relaciones bajas o no significativas con la mayoría de los factores (r = .01 – .11), lo que sugiere que opera de manera relativamente independiente dentro del modelo de percepción. Este comportamiento indica que la percepción de riesgo no se integra de forma directa con las valoraciones positivas de utilidad, confianza o autonomía, sino que constituye un componente con dinámica propia dentro de la estructura general del instrumento. Por su parte, las políticas institucionales evidenciaron asociaciones moderadas con los factores positivos (r = .25 – .30, p < .01), indicando que quienes valoran más la herramienta reconocen también la importancia de contar con lineamientos claros para su uso académico.

Estos patrones respaldan el carácter multidimensional del instrumento, mostrando que la valoración positiva de ChatGPT no implica necesariamente una disminución en la percepción de riesgos, sino que ambos constructos se comportan de manera diferenciada. En otras palabras, los estudiantes pueden identificar ventajas académicas de la herramienta y, al mismo tiempo, mantener reservas respecto de sus implicaciones éticas, normativas o de uso inapropiado.

\subsection{Análisis de diferencias por género}

Los análisis de varianza para comparar a hombres y mujeres en las dimensiones evaluadas se presentan en la Tabla \ref{tab3}. No se observaron diferencias estadísticamente significativas entre géneros en ninguna de las dimensiones evaluadas, incluidas la utilidad percibida, la autonomía en el aprendizaje, los riesgos percibidos, las competencias cognitivas, la confianza en las respuestas proporcionadas por la herramienta y la percepción de las políticas institucionales (p > .05 en todos los casos). Estos resultados indican que, en la muestra estudiada, las valoraciones generales sobre el uso académico de ChatGPT fueron similares entre mujeres y hombres en casi todas las dimensiones analizadas.

\begin{table}[htbp]
\centering
\begin{threeparttable}
\caption{Resultados del ANOVA por género.}
\label{tab3}
\begin{tabular}{lcc}
\toprule
Dimensión & F & p (Sig.) \\ 
\midrule
UTI & 0.006 & .937 \\
AUT & 0.077 & .782 \\
RIE & 0.079 & .779 \\
COM & 0.020 & .888 \\
CON & 0.148 & .701 \\
POL & 0.742 & .390 \\
CRI & 2.763 & .097 \\
\bottomrule
\end{tabular}
\source{Estadísticos calculados por los autores en SPSS con datos del estudio.}
\end{threeparttable}
\end{table}

En la dimensión de criterio ético se observó una diferencia de mayor magnitud relativa que en el resto de las dimensiones (p = .097); sin embargo, esta no alcanzó significancia estadística. Por ello, aunque este resultado sugiere cierta variación en la forma en que ambos grupos valoran los aspectos éticos asociados al uso de la herramienta, no es suficiente para afirmar una diferencia concluyente en términos estadísticos. En conjunto, los hallazgos muestran que la percepción y el uso académico de ChatGPT no presentan diferencias estadísticamente significativas entre mujeres y hombres en las dimensiones evaluadas.

\subsection{Diferencias entre nivel de preparación académica}

La Tabla \ref{tab4} presenta los resultados del ANOVA realizado para identificar diferencias entre los niveles de preparación académica (básico, intermedio y avanzado). Únicamente la dimensión de autonomía en el aprendizaje mostró diferencias estadísticamente significativas entre grupos (p < .01). Las demás dimensiones no presentaron variaciones estadísticamente significativas (p > .05), aunque utilidad percibida y políticas institucionales registraron valores p menores que el resto de las dimensiones analizadas, sin alcanzar el criterio establecido para significancia.

\begin{table}[htbp]
\centering
\begin{threeparttable}
\caption{Resultados del ANOVA por nivel académico.}
\label{tab4}
\begin{tabular}{lcc}
\toprule
Dimensión & F & p (Sig.) \\ 
\midrule
UTI & 2.864 & .058 \\
AUT & 5.198 & .006 \\
RIE & 0.480 & .619 \\
COM & 1.998 & .137 \\
CON & 1.245 & .289 \\
POL & 2.484 & .085 \\
CRI & 0.725 & .485 \\
\bottomrule
\end{tabular}
\vspace{6pt}
\source{Estadísticos calculados por los autores en SPSS con datos del estudio.}
\end{threeparttable}
\end{table}

En el caso de autonomía, los análisis post hoc (Tukey) indicaron que el nivel intermedio presenta una percepción significativamente mayor que el nivel básico (p = .011). El nivel avanzado también mostró medias superiores, aunque sin alcanzar significancia estadística en las comparaciones múltiples. Este patrón sugiere que la percepción de autonomía asociada al uso de ChatGPT aumenta conforme avanzan los estudios, particularmente entre etapas intermedias y avanzadas de la formación universitaria. No obstante, esta tendencia debe interpretarse con cautela, dado que el grupo avanzado estuvo conformado por un número considerablemente menor de participantes que los otros niveles.

En conjunto, los resultados evidencian que el nivel académico influye de manera específica en la percepción de autonomía, mientras que las demás dimensiones evaluadas —utilidad percibida, riesgos, competencias cognitivas, confianza, políticas institucionales y criterio ético— no mostraron diferencias estadísticamente significativas entre los distintos niveles de formación. Esto sugiere que la trayectoria académica no modifica de manera general todas las percepciones sobre ChatGPT, sino que su efecto se concentra de forma más clara en la experiencia subjetiva de autonomía en el aprendizaje.

\subsection{Diferencias entre categoría de edad}

La Tabla \ref{tab5} presenta los resultados del análisis de varianza aplicado a las tres categorías de edad consideradas en el estudio (jóvenes, adultos emergentes y adultos profesionales). Los análisis no identificaron diferencias estadísticamente significativas en ninguna de las dimensiones evaluadas (p > .05). En términos generales, los resultados muestran que la edad no se asoció con cambios sistemáticos en las percepciones estudiantiles sobre el uso académico de ChatGPT dentro de la muestra analizada.

\begin{table}[htbp]
\centering
\begin{threeparttable}
\caption{Resultados del ANOVA por categoría de edad.}
\label{tab5}
\begin{tabular}{lcc}
\toprule
Dimensión & F & p (Sig.) \\ 
\midrule
UTI & 0.540 & .583 \\
AUT & 0.595 & .552 \\
RIE & 1.802 & .166 \\
COM & 2.323 & .099 \\
CON & 0.141 & .869 \\
POL & 0.815 & .443 \\
CRI & 0.231 & .794 \\
\bottomrule
\end{tabular}
\vspace{6pt}
\source{Estadísticos calculados por los autores en SPSS con datos del estudio.}
\end{threeparttable}
\end{table}

Las pruebas post hoc mostraron subconjuntos homogéneos, lo que indica que las percepciones sobre utilidad, autonomía, riesgos, competencias cognitivas, confianza, políticas institucionales y criterio ético no presentaron diferencias estadísticamente significativas entre los distintos rangos de edad representados en la muestra. De esta manera, los datos sugieren que las valoraciones sobre ChatGPT se mantuvieron relativamente semejantes entre estudiantes jóvenes y de mayor edad, al menos en las dimensiones consideradas por el instrumento.

\section{Discusión}

Los resultados de este estudio aportan evidencia relevante sobre cómo el estudiantado percibe el uso académico de ChatGPT en un contexto universitario específico. En términos generales, las percepciones fueron moderadamente favorables, lo cual resulta consistente con investigaciones recientes que destacan el atractivo pedagógico de estas tecnologías por su rapidez, accesibilidad y capacidad de apoyo en tareas de lectura, escritura, resolución de problemas y consulta académica \cite{Stone2025,SunLiHossainZabin2025}. No obstante, la variabilidad observada en algunos reactivos muestra que estas valoraciones no fueron completamente homogéneas, lo que sugiere la presencia de matices en la forma en que el estudiantado interpreta la utilidad y el alcance de la herramienta.

Un hallazgo relevante del estudio es la ausencia de diferencias estadísticamente significativas por género y categoría de edad en ninguna de las dimensiones evaluadas. Este resultado coincide con trabajos recientes que señalan que las actitudes hacia ChatGPT tienden a mantenerse relativamente estables entre distintos grupos demográficos, dado que su adopción suele vincularse más con la utilidad percibida que con diferencias socioculturales amplias \cite{Haq2025,Pitts2025}. Asimismo, estudios globales han planteado que la familiaridad digital compartida entre mujeres y hombres puede reducir brechas tradicionales en la apropiación tecnológica dentro de la educación superior \cite{MartinMoncunillAlonsoMartinez2025}. En este sentido, los datos obtenidos sugieren que, al menos en la muestra analizada, ni el género ni la edad operaron como factores diferenciadores relevantes en la percepción del uso académico de ChatGPT.


En contraste, sí se encontraron diferencias significativas en la percepción de autonomía según el nivel de preparación académica. De manera específica, el nivel intermedio reportó una percepción significativamente mayor que el nivel básico, mientras que el nivel avanzado mostró medias superiores, aunque sin alcanzar significancia estadística en las comparaciones múltiples. Estos resultados pueden interpretarse a la luz de investigaciones que sugieren que la experiencia acumulada en la trayectoria universitaria favorece estrategias más elaboradas de consulta, verificación y aprovechamiento de herramientas de IA, fortaleciendo con ello procesos de autorregulación y confianza en el aprendizaje asistido por IA \cite{Dong2024,Hong2025}. No obstante, el comportamiento del grupo avanzado debe interpretarse con cautela debido al tamaño reducido de la muestra en ese nivel.

En cuanto a los riesgos percibidos, los resultados muestran correlaciones bajas o no significativas con el resto de las dimensiones, lo que refuerza estudios que identifican esta dimensión como un componente relativamente diferenciado dentro de la percepción del uso académico de la IA, asociado a preocupaciones sobre privacidad, veracidad de la información, dependencia cognitiva y uso inapropiado en procesos evaluativos \cite{Bender2021,Jose2025}. En este punto, los hallazgos del estudio sugieren que el reconocimiento de beneficios académicos no elimina, por sí mismo, la presencia de reservas frente a los posibles efectos problemáticos de la herramienta. Estas preocupaciones han llevado a un creciente llamado internacional para revisar modelos de evaluación y reforzar políticas institucionales de integridad académica \cite{Cotton2024,Marin2025,SolSokHeng2025,WangCornely2023}.

Algunos de estos patrones también pueden relacionarse con debates emergentes sobre la externalización de procesos cognitivos en el uso de modelos de lenguaje. Desde esta perspectiva, el uso frecuente de ChatGPT puede favorecer determinadas tareas académicas, pero también generar reservas si no existe una mediación pedagógica adecuada que oriente su empleo dentro de procesos formativos estructurados \cite{Hong2025,Jose2025}. Esto guarda relación con el patrón observado en el presente estudio: las valoraciones positivas sobre utilidad y autonomía coexistieron con una percepción de riesgos que se comportó de manera relativamente diferenciada frente al resto de las dimensiones.

Finalmente, los resultados reflejan la necesidad de fortalecer lineamientos institucionales que orienten el uso responsable de IA en entornos educativos. Como sugieren \textcite{GarciaLopezTrujilloLinan2025}, la regulación debe equilibrar innovación pedagógica e integridad académica, integrando principios de transparencia, ética y alfabetización digital. En el caso del presente estudio, esta necesidad encuentra respaldo empírico en las puntuaciones más bajas observadas en los reactivos vinculados con políticas institucionales, lo que sugiere una percepción limitada sobre la claridad normativa existente. En conjunto, los hallazgos muestran que la percepción estudiantil sobre ChatGPT combina aceptación funcional, diferencias acotadas según la trayectoria académica y reservas específicas en torno a riesgos y regulación, aportando evidencia útil para el diseño de lineamientos institucionales más ajustados a la experiencia real del estudiantado.

\section{Conclusiones}

El presente estudio contribuye al entendimiento de cómo estudiantes universitarios perciben el uso académico de ChatGPT, ofreciendo evidencia empírica que complementa discusiones teóricas recientes sobre IAG en educación superior. Los resultados muestran una valoración general moderadamente favorable del uso académico de la herramienta, con especial reconocimiento de su utilidad y con asociaciones positivas entre utilidad, autonomía y competencias cognitivas. Al mismo tiempo, se identificaron áreas sensibles relacionadas con los riesgos percibidos, el criterio ético y las políticas institucionales.

Una de las aportaciones centrales del estudio es la identificación de diferencias significativas únicamente en la dimensión de autonomía según nivel académico. En particular, el nivel intermedio presentó una percepción de autonomía mayor que el nivel básico, mientras que el nivel avanzado mostró medias superiores sin alcanzar significancia estadística. Esto sugiere que, en la muestra analizada, las diferencias observadas se asociaron más con la trayectoria académica que con el género o la edad. En línea con investigaciones recientes \cite{Haq2025,Pitts2025}, la experiencia universitaria parece desempeñar un papel relevante en la forma en que algunos estudiantes integran la IA en sus procesos de aprendizaje.

Asimismo, los hallazgos sobre riesgos percibidos y preocupaciones éticas, junto con las puntuaciones más bajas en los reactivos vinculados con políticas institucionales, resaltan la necesidad de desarrollar marcos institucionales más claros que orienten el uso responsable de la IA en contextos educativos. La literatura indica que instituciones de educación superior en todo el mundo están enfrentando desafíos similares, desde el aumento de prácticas deshonestas hasta la redefinición de la originalidad académica \cite{Abdurrahim2025,Chen2024,Cotton2024,VanNiekerkDelportSutherland2025}. Por ello, los resultados de este estudio pueden servir como insumo para fortalecer políticas locales, talleres de alfabetización digital y estrategias de evaluación adaptadas a la era de la IAG.

Entre las limitaciones del estudio se encuentran la naturaleza transversal de los datos, el uso de una muestra no probabilística por conveniencia, la realización del estudio en una sola institución y la baja representación de estudiantes en niveles avanzados, lo cual sugiere que futuras investigaciones podrían beneficiarse de muestreos más equilibrados, diseños longitudinales y análisis cualitativos que profundicen en prácticas de uso real. También sería pertinente explorar con mayor detalle cómo los modelos de lenguaje influyen en la construcción de conocimiento, la autoría académica y el desarrollo del pensamiento crítico.

En síntesis, los hallazgos evidencian que ChatGPT es percibido como un recurso académico relevante por el estudiantado, con ventajas percibidas en términos de utilidad y apoyo al aprendizaje, pero también con desafíos éticos y regulatorios que requieren atención institucional. Un enfoque equilibrado que combine innovación, integridad y alfabetización digital será indispensable para orientar el uso responsable de la IAG en la educación superior contemporánea.



\printbibliography\label{sec-bib}
% if the text is not in Portuguese, it might be necessary to use the code below instead to print the correct ABNT abbreviations [s.n.], [s.l.]
%\begin{portuguese}
%\printbibliography[title={Bibliography}]
%\end{portuguese}


%full list: conceptualization,datacuration,formalanalysis,funding,investigation,methodology,projadm,resources,software,supervision,validation,visualization,writing,review
\begin{contributors}[sec-contributors]
\authorcontribution{Gloria Mónica Martínez Aguilar}[conceptualization,datacuration,formalanalysis,investigation,projadm,methodology,supervision,writing,review]
\authorcontribution{Kristian Segura Félix}[datacuration,investigation,visualization,review]
\authorcontribution{Geshel David Guerrero López}[datacuration,investigation,visualization,review]
\end{contributors}


\begin{dataavailability}
\txtdataavailability{dataonly} % options: dataavailable, dataonly, databody, datanotav, nodata
\end{dataavailability}

\section*{\fontsize{10}{12}\selectfont Declaración ética e inteligencia artificial}
La investigación fue realizada respetando los principios éticos aplicables a estudios en contextos educativos, garantizando la participación voluntaria, el anonimato y la confidencialidad de los datos recolectados.

Asimismo, se declara que herramientas de inteligencia artificial generativa (ChatGPT, OpenAI) fueron utilizadas únicamente como apoyo en procesos de redacción, revisión lingüística y organización textual del manuscrito. El análisis, la interpretación de los datos y la versión final del artículo son responsabilidad exclusiva de los autores, quienes garantizan la originalidad y la integridad académica del trabajo presentado.

\end{document}

